296 5 Differential Calculus

whose solution $f(x)$ is uniquely determined when the ``initial condition'' $f(0)$ is
specified. Then
$$f(x) = f(0)e^{\alpha x}.$$
We introduced the number $e$ and the function $e^x = \exp x$ in a rather formal way
earlier, assuring the reader that $e$ really was an important number and $\exp x$ really
was an important function. It is now clear that even if we had not introduced this
function earlier, it would certainly have been necessary to introduce it as the solution
of the important, though very simple equation (5.152). More precisely, it would have
sufficed to introduce the function that is the solution of Eq. (5.152) for some specific
value of $\alpha$, for example, $\alpha = 1$; for the general equation (5.152) can be reduced to
this case by changing to a new variable $t$ connected with $x$ by the relation $x = \frac{t}{\alpha}$
$(\alpha \ne 0)$.
Indeed, we then have
$$f(x) = f\left(\frac{t}{\alpha}\right) = F(t), \quad \frac{df(x)}{dx} = \frac{dF(t)}{dt} \frac{dt}{dx} = \alpha F'(t)$$
and instead of the equation $f'(x) = \alpha f(x)$ we then have $\alpha F'(t) = \alpha F(t)$, or $F'(t) =
F(t)$.
Thus, let us consider the equation
$$f'(x) = f(x) \quad (5.153)$$
and denote the solution of this equation satisfying $f(0) = 1$ by $\exp x$.
Let us check to see whether this definition agrees with our previous definition of
$\exp x$.
Let us try to calculate the value of $f(x)$ starting from the relation $f(0) = 1$ and
the assumption that $f$ satisfies (5.153). Since $f$ is differentiable, it is continuous.
But then Eq. (5.153) implies that $f'(x)$ is also continuous. Moreover, it follows
from (5.153) that $f$ also has a second derivative $f''(x) = f'(x)$, and in general that
$f$ is infinitely differentiable. Since the rate of variation $f'(x)$ of the function $f(x)$
is continuous, the function $f'$ changes very little over a small interval $h$ of variation
of its argument. Therefore $f(x_0 + h) = f(x_0) + f'(\xi)h \approx f(x_0) + f'(x_0)h$. Let us
use this approximate formula and traverse the interval from $0$ to $x$ in small steps of
size $h = \frac{x}{n}$, where $n \in \mathbb{N}$. If $x_0 = 0$ and $x_{k+1} = x_k + h$, we should have
$$f(x_{k+1}) \approx f(x_k) + f'(x_k)h.$$
Taking account of (5.153) and the condition $f(0) = 1$, we have
$$f(x) = f(x_n) \approx f(x_{n-1}) + f'(x_{n-1})h = f(x_{n-1})(1 + h) \sim (f(x_{n-2}) + f'(x_{n-2})h) (1 + h) =$$